\documentclass[11pt,a4paper]{article}
\usepackage[margin=22mm]{geometry}
\usepackage{amsmath,amssymb,booktabs,graphicx,microtype,hyperref,array}
\allowdisplaybreaks
\title{Temporary fiber probes and spatial mismatch\\A carried-coefficient forward derivation}
\author{Illustrative spatial scenarios; no fit to final contrast}
\date{}
\begin{document}\maketitle

\noindent\textbf{What this test answers.} A temporary fiber probe can constrain
spatial modes, but its coupling efficiency is not attenuation in the operating
interferometer after the probe is removed. We test two explicit spatial
arrangements that give the same example probe efficiencies, @@ETA_A_PERCENT@@\% for A and
@@ETA_B_PERCENT@@\% for B. We then apply the same illustrative C/D angular mismatch to both
arrangements. One yields modest contrast and the other high contrast. Thus a
spatial mechanism can produce high contrast, but the two efficiencies do not
identify which mechanism the apparatus has. These efficiency values are examples, not
newly acquired data; the geometric assumptions and tilt are not fitted.

\section{What is measured, assumed, and still unknown}
The user reports temporary steering mirrors into a fiber z-fold. A/B share one
probe fiber; C/D share another. The probes are removed from the operating path.
It is not yet confirmed that each fiber position, focusing relay and calibration
remain fixed while switching the two paths. All overlap inferences below are
conditional on that common reference. Independently optimizing each path into
the fiber does not measure overlap with one fixed reference mode.

\begin{center}\small\begin{tabular}{>{\raggedright\arraybackslash}p{.43\linewidth}>{\raggedright\arraybackslash}p{.50\linewidth}}\toprule
Quantity & Role in this calculation \\\midrule
NPBS power ratios & Earlier direct-meter data; C/D completions kept separate \\
Probe efficiencies $\eta_A=@@ETA_A@@$, $\eta_B=@@ETA_B@@$ & Illustrative input examples \\
Common probe throughput $\eta_0=@@ETA_0@@$ & Assumed, not measured \\
Equal Gaussian widths and curvature plane & Explicit geometric approximation \\
Signs of A/B transverse displacement & Two alternative hypotheses \\
D-arm slope $q=@@Q@@$ & Chosen example; additionally swept, not fitted \\
Beam radius $w$, wavelength $\lambda$ & Optional conversion inputs; not measured by probe powers \\
Input amplitudes $a_x=a_y=1$ & Retained equal-input baseline \\
\bottomrule\end{tabular}\end{center}
@@PHYSICAL_GEOMETRY@@

No probe loss, new retarder, polarization mixing or fitted phase offset is added.
Final detection in this test integrates the full free-space beam in each output.
Finite apertures, a permanent fiber receiver, or polarization-selective detection
are different observables.

\section{What a fiber coupling efficiency actually constrains}
Normalize the transverse mode $u_j$ of path $j$ and the probe fiber mode $f$:
$\langle u_j,u_j\rangle=\langle f,f\rangle=1$. The complex field coupling is
\begin{equation}
h_j=\langle f,u_j\rangle=\iint f^*(\mathbf r)u_j(\mathbf r)\,d^2r,
\qquad \eta_j=\eta_0|h_j|^2.
\end{equation}
For the configured common throughput, the example determines
$|h_A|=\sqrt{\eta_A/\eta_0}=@@HA@@$ and $|h_B|=\sqrt{\eta_B/\eta_0}=@@HB@@$.
It does not determine their phases or the mutual free-space overlap
$\Gamma_{AB}=\langle u_A,u_B\rangle$. To see why, decompose
\[
u_j=h_j f+v_j,\qquad\langle f,v_j\rangle=0,\qquad
\|v_j\|^2=1-|h_j|^2.
\]
Then
\begin{equation}
\Gamma_{AB}=h_A^*h_B+\langle v_A,v_B\rangle,\qquad
|\Gamma_{AB}-h_A^*h_B|\le\sqrt{(1-|h_A|^2)(1-|h_B|^2)}.
\end{equation}
For the configured values the disk center has magnitude @@GRAM_CENTER@@ and
radius @@GRAM_RADIUS@@. The allowed magnitude interval is
$@@GRAM_LOW@@\le|\Gamma_{AB}|\le@@GRAM_HIGH@@$.
This broad bound allows very different modes. It is not a measured visibility.
Common lens/Fresnel/instrument losses must be separated from true mode coupling;
otherwise even the quoted geometrical bound is conditional on $\eta_0$.

\section{A simple geometric model: displacement and angular tilt}
To construct physically consistent examples, use normalized equal-waist modes
at one common plane:
\begin{equation}
u(\mathbf r;\mathbf d,\boldsymbol\kappa)=\sqrt{\frac{2}{\pi w^2}}
\exp\left[-\frac{|\mathbf r-\mathbf d|^2}{w^2}\right]
\exp(i\boldsymbol\kappa\cdot\mathbf r).
\end{equation}
Here $w$ is the $1/e$ field radius ($1/e^2$ intensity radius). A small angular
tilt has $\boldsymbol\kappa\simeq k\boldsymbol\theta$, $k=2\pi/\lambda$.
Introduce dimensionless $\mathbf D=\mathbf d/w$ and $\mathbf Q=w\boldsymbol\kappa$.
Completing the square in the Gaussian overlap integral gives
\begin{equation}
\boxed{G_{ab}=\langle u_a,u_b\rangle=
\exp\left[-\frac{|\mathbf D_a-\mathbf D_b|^2}{2}
-\frac{|\mathbf Q_a-\mathbf Q_b|^2}{8}
+\frac{i}{2}(\mathbf Q_b-\mathbf Q_a)\cdot(\mathbf D_a+\mathbf D_b)\right].}
\end{equation}
For the centered, untilted probe mode, this reduces to
\begin{equation}
h(\mathbf D,\mathbf Q)=\exp[-|\mathbf D|^2/2-|\mathbf Q|^2/8+i\mathbf Q\cdot\mathbf D/2],
\quad\boxed{\eta=\eta_0\exp[-|\mathbf D|^2-|\mathbf Q|^2/4].}
\end{equation}
One efficiency therefore constrains a combination of displacement and tilt,
not either individually. If \emph{only displacement} is assumed, then
\[
D_A=\sqrt{\ln(\eta_0/\eta_A)}=@@DA@@,\qquad
|D_B|=\sqrt{\ln(\eta_0/\eta_B)}=@@DB@@.
\]
The two hypotheses are $D_B=+@@DB@@$ (same side of the probe mode) and
$D_B=-@@DB@@$ (opposite sides). Their A/B separations are respectively
$@@SEP_SAME@@\,w$ and $@@SEP_OPPOSITE@@\,w$. Both give exactly the same probe
power efficiencies. If tilt instead caused all the loss, it would satisfy
$|\theta|=2\sqrt{\ln(\eta_0/\eta)}/(kw)$; no numerical angle can be recovered
without $w$, $\lambda$, and the geometric assumptions.
\begin{center}\includegraphics[width=\linewidth]{geometry.pdf}\end{center}

\section{Propagate four route modes, without inserting probe loss}
Keep the existing splitter fields and coefficients. Define
\[
\alpha=t_{1x}t_{2x},\quad\beta=r_{1x}r_{2x},\quad
\gamma=r_{1y}t_{2y},\quad\delta=t_{1y}r_{2y},\quad z=e^{i\phi_2}.
\]
The splitter field coefficients are carried from $t=\sqrt{P_t/(P_t+P_r)}$,
$r=\sqrt{P_r/(P_t+P_r)}$ (all raw powers in $\mu$W):
\begin{center}\begin{tabular}{lrrr}\toprule
Input & $\mathcal T=P_t/(P_t+P_r)$ & $t$ & $r$ \\\midrule
@@SPLITTER_ROWS@@
\bottomrule\end{tabular}\end{center}
The C-reference completion uses the C rows for both reciprocal NPBS2 inputs;
the D-reference completion uses the D rows instead. They are separate lossless
completions, not an average or fit. Reflection retains the ideal factor $i$.
The ideal PBS sends x through A and y through B. Each route now carries its
own normalized spatial mode rather than an assumed identical mode:
\begin{align}
E_x(\mathbf r)&=a_xe^{i\phi_1}(\alpha z u_{AC}-\beta u_{AD}),\\
E_y(\mathbf r)&=ia_y(\gamma z u_{BC}+\delta u_{BD}).
\end{align}
The relative angular mismatch is implemented as a spatial phase ramp in D,
$u_{AD}=e^{i\kappa x}u_A$, $u_{BD}=e^{i\kappa x}u_B$, while
$u_{AC}=u_A$, $u_{BC}=u_B$. This is a unitary transformation of transverse
fields: it redistributes spatial phase and removes no power. It represents
a relative C/D angular mismatch at the common output reference plane, not a
known measured mirror setting or a specified z-fold geometry.

A full-aperture power meter integrates $|E_x|^2+|E_y|^2$. Define
$\Gamma_A=\langle u_{AD},u_{AC}\rangle$ and
$\Gamma_B=\langle u_{BD},u_{BC}\rangle$. Explicit expansion gives
\begin{align}
I_{E,x}&=a_x^2[\alpha^2+\beta^2-2\alpha\beta\Re(z\Gamma_A)],\\
I_{E,y}&=a_y^2[\gamma^2+\delta^2+2\gamma\delta\Re(z\Gamma_B)].
\end{align}
Thus, carrying independent input amplitudes,
\begin{align}
A_E&=a_x^2(\alpha^2+\beta^2)+a_y^2(\gamma^2+\delta^2),\\
K_E&=a_y^2\gamma\delta\Gamma_B-a_x^2\alpha\beta\Gamma_A,\\
\boxed{I_E}&=\boxed{A_E+2\Re(K_Ee^{i\phi_2})
=A_E+2\Re K_E\cos\phi_2-2\Im K_E\sin\phi_2.}
\end{align}
The phase $\phi_1$ still cannot modulate total intensity: x and y are orthogonal.
Scalar spatial mismatch alone in this restricted optical network does not
explain a phi1 power fringe. It can alter the polarization and degree of
polarization through spatial integration.

For the measured C-reference splitter coefficients from the prior report,
\[
\alpha\beta=@@X_WEIGHT@@,\quad\gamma\delta=@@Y_WEIGHT@@,\quad
A_E\big|_{a_x=a_y=1}=@@MEAN@@,
\]
so the spatial extension is exactly
\begin{equation}
K_E=@@Y_WEIGHT@@\,\Gamma_B-@@X_WEIGHT@@\,\Gamma_A.
\end{equation}
No old fitted retarder or fitted contrast coefficient appears here.

For F let $h=t_{1x}r_{2x}$, $j=r_{1x}t_{2x}$,
$k=t_{1y}t_{2y}$ and $\ell=r_{1y}r_{2y}$. Then
\begin{align}
F_x&=ia_xe^{i\phi_1}(hz u_{AC}+ju_{AD}),&
F_y&=a_y(ku_{BD}-\ell z u_{BC}),\\
I_F&=a_x^2[h^2+j^2+2hj\Re(z\Gamma_A)]
+a_y^2[k^2+\ell^2-2k\ell\Re(z\Gamma_B)].
\end{align}
Since $hj=\alpha\beta$, $k\ell=\gamma\delta$ and all spatial modes have unit
norm, $I_E+I_F=a_x^2+a_y^2=2$. This assumes complete collection and the same
mode coordinate convention across the lossless final splitter.

\section{Derive contrast and show how high contrast can appear}
Write $K_E=|K_E|e^{i\psi}$. Then
$I_E=A_E+2|K_E|\cos(\phi_2+\psi)$ and
\begin{equation}
\boxed{V_E=\frac{I_{\max}-I_{\min}}{I_{\max}+I_{\min}}
=\frac{2|K_E|}{A_E}},\qquad V_F=\frac{2|K_E|}{2-A_E}.
\end{equation}
For the ideal equal-input splitter baseline, $A_E=1$,
$\alpha\beta=\gamma\delta=1/4$, so
\begin{equation}
\boxed{V_E=\frac{|\Gamma_B-\Gamma_A|}{2}.}
\end{equation}
\textbf{A shared visibility reduction is insufficient.} If
$\Gamma_A=\Gamma_B$, the two polarization fringes still cancel. If both overlaps
are merely nonnegative real numbers, their difference cannot exceed one, so
$V_E\le1/2$ in the ideal equal-input case. High contrast such as 0.6 requires
more than a common real overlap reduction in this model.

For the explicit Gaussian D-arm ramp with $q=\kappa w$,
\[
\Gamma_A=e^{-q^2/8}e^{-iqD_A},\qquad
\Gamma_B=e^{-q^2/8}e^{-iqD_B}.
\]
Using $|e^{ia}-e^{ib}|=2|\sin[(a-b)/2]|$ gives the carried geometric result
\begin{equation}
\boxed{V_E(q)=e^{-q^2/8}
\left|\sin\!\left(\frac{q(D_A-D_B)}{2}\right)\right|}
\quad\text{(ideal splitters, equal input).}
\end{equation}
The ramp reduces individual overlap magnitudes but also gives the A-derived
and B-derived interference terms different phases because their beam centers
differ. Their formerly cancelling fringes need not cancel. No polarization
rotation is required for this particular spatial mechanism.

At the \emph{chosen, unmeasured} example $q=@@Q@@$, the common overlap magnitude
is $e^{-q^2/8}=@@GFACT@@$. Substitution into the measured-splitter expression gives:
@@NUMERICS@@
The scenario is not selected by optimizing agreement with $V=0.6$. The sweep
below exposes the full response rather than fitting the tilt.
\begin{center}\small\begin{tabular}{llrrr}\toprule
Geometry & Splitter case & $V_E$ & $100V_E$ (\%) & $V_F$ \\\midrule
@@RESULT_ROWS@@
\bottomrule\end{tabular}\end{center}
\begin{center}\includegraphics[width=\linewidth]{visibility.pdf}\end{center}

\section{Carry the spatial overlap into Stokes, not just intensity}
For integrated detection define the coherency matrix
$J_{pq}=\iint E_p(\mathbf r)E_q^*(\mathbf r)\,d^2r$.
Then, with the repository convention,
\[
S_0=J_{xx}+J_{yy},\quad S_1=J_{yy}-J_{xx},\quad
S_2=2\Re J_{xy},\quad S_3=2\Im J_{xy},\quad s_n=S_n/S_0.
\]
The diagonal entries are the $I_{E,x}$ and $I_{E,y}$ already derived. Expanding
the cross term retains all four spatial overlaps:
\begin{align}
J_{xy,E}=-ia_xa_ye^{i\phi_1}\big[&
\alpha\gamma\langle u_{BC},u_{AC}\rangle
+\alpha\delta z\langle u_{BD},u_{AC}\rangle\\
&-\beta\gamma z^*\langle u_{BC},u_{AD}\rangle
-\beta\delta\langle u_{BD},u_{AD}\rangle\big].
\end{align}
For F the corresponding expression is
\begin{align}
J_{xy,F}=ia_xa_ye^{i\phi_1}\big[&hkz\langle u_{BD},u_{AC}\rangle
-h\ell\langle u_{BC},u_{AC}\rangle\\
&+jk\langle u_{BD},u_{AD}\rangle
-j\ell z^*\langle u_{BC},u_{AD}\rangle\big].
\end{align}
Insert the explicit Gaussian $G_{ab}$ into these products, then take real and
imaginary parts and divide by the actual detected $S_0$. These are the plotted
Stokes parameters; no per-curve scaling or forced unit Stokes norm is applied.
The degree of polarization is
\[
\mathcal P=\sqrt{S_1^2+S_2^2+S_3^2}/S_0\le1.
\]
Different polarizations can occupy different spatial modes, so full-plane
integration can yield $\mathcal P<1$ even for a deterministic coherent laser
field. A single pure Jones vector then does not describe the integrated state.
When $\mathcal P>0$, the direction $\mathbf s/\mathcal P$ describes the polarized
component; its ellipse angles can be obtained from
$\theta=\tfrac12\operatorname{atan2}(s_2,s_1)$ and
$\epsilon=\tfrac12\arcsin(s_3/\mathcal P)$ in this convention. At
$\mathcal P=0$, no unique ellipse exists.

\section{What the second probe fiber can and cannot establish}
Apply the same overlap formula to C/D at its own calibrated probe plane and
its own fixed receiver mode. Do not multiply efficiencies from two temporary
probe stations together as though both fibers stayed in the optical train.
The relay mapping between the two stations also has to be characterized before
identifying a common physical displacement or angular mismatch.

With both A and B illuminated, C/D generally contain two polarization-dependent
spatial profiles. To constrain the overlaps used above, block down to the four
routes AC, AD, BC, BD and measure each at one unchanged probe alignment, with a
free-space reference power for that same route. For our illustrative Gaussian
case the geometric probe efficiencies would be:
\begin{center}\begin{tabular}{lrr}\toprule
Route & Same-side geometry & Opposite-side geometry \\\midrule
@@PROBE_ROWS@@
\bottomrule\end{tabular}\end{center}
They are identical in the two geometries even though their free-space contrasts
differ strongly: scalar probe powers do not supply the displacement signs or
complex overlap phases. These are hypothetical readings at the common modeled
plane, not predictions for an uncharacterized second z-fold relay.

Useful next acquisitions are isolated-route free-space/fiber power ratios with
uncertainty and a fixed fiber reference, coupler x/y scans with signed coordinates,
and camera centroids at two longitudinal planes. The latter separate displacement
from angle. Record width convention, relay settings and power calibration.
Dedicated coherent overlap measurements can further constrain phase. If C/D
probe data disagree with these illustrative profiles, replace the geometric
hypothesis using those independent observations; do not tune it to final contrast.

\section{Different question: what if a fiber really stays in the path?}
A permanent single-mode receiver measures the projected Jones field
$E_{p,f}=\sum_a h_a e_{a,p}$, not the free-space integral above. If its coefficients
only distinguish A-origin versus B-origin light, $h_{AC}=h_{AD}=h_A$ and
$h_{BC}=h_{BD}=h_B$, then for the ideal hybrid network
\[
I_{E,f}=\eta_A\sin^2(\phi_2/2)+\eta_B\cos^2(\phi_2/2),\qquad
V=\frac{|\eta_A-\eta_B|}{\eta_A+\eta_B}=@@INLINE_V@@.
\]
That is a @@INLINE_PERCENT@@\% example, and is \emph{not} the correct way to treat the removed
alignment probes. In an ordinary same-polarization two-arm MZI the corresponding
balanced-recombination expression instead gives
$V=2\sqrt{\eta_A\eta_B}/(\eta_A+\eta_B)=@@ORDINARY_V@@$ under perfect coherence.
Orthogonal polarization, spatial overlap and the detector observable determine
which formula applies. Coupling imbalance by itself is not a universal visibility.

\section{Reference and checks}
The overlap-integral definition and Gaussian mode-matching framework are
consistent with Schwab et al., \emph{Optics Express} 30, 32292 (2022),
\href{https://doi.org/10.1364/OE.465101}{doi:10.1364/OE.465101}, Sec.~2.3.
The equal-waist overlap and carried-coefficient interferometer formulas above
are evaluated explicitly here. Tests compare analytical overlaps with numerical
integration, harmonic visibility with direct propagation, total power conservation,
probe ambiguity and partial-polarization Stokes. Figures use optical phase in
degrees; $q$ is a dimensionless spatial phase slope, not another swept arm phase.
All retained spatial numbers are labeled hypotheses until acquired independently.

\clearpage\begin{center}\includegraphics[width=\linewidth,height=.92\textheight,keepaspectratio]{phase-E.pdf}\end{center}
\clearpage\begin{center}\includegraphics[width=\linewidth,height=.92\textheight,keepaspectratio]{phase-F.pdf}\end{center}
\end{document}
